% ============================================================
% Chapter 6: Conclusions and Recommendations
% File: chapters/conclusions.tex
% ============================================================

\section{Conclusions}
\label{sec:conclusions}

This thesis set out to evaluate, under a common and reproducible
protocol, whether the choice of functional basis materially affects the
representation of volatile economic time series and the functional
principal component analysis computed from it. The central conclusion is
that it does, and that the effect cannot be read from reconstruction
error alone. On the extended panel of annual GDP growth rates for 20
Latin American countries (1960--2024), the B-spline, Fourier, and
wavelet families reach comparable mean reconstruction errors (2.51,
2.65, and 2.52--2.53 percentage points, respectively;
Table~\ref{tab:rmse_summary}), yet they differ in which features of the
trajectories they retain, and these differences propagate directly into
the number, shape, and economic interpretation of the estimated
functional principal components. The basis expansion is therefore not a
neutral preprocessing step: it is a modelling decision that determines
which structure of the data is available to every subsequent analysis.

The specific conclusions are organized below according to the specific
objectives stated in Section~\ref{sec:objectives}.

\subsection{Extension of the dataset}

The annual GDP growth panel of \cite{padilla2020} was extended to
1960--2024 for 20 Latin American countries using World Bank Open Data,
incorporating the 2020 COVID-19 contraction and the subsequent recovery.
The extended dataset preserves the variance structure of the original
study: the FPCA of the unsmoothed data yields a cumulative variance of
47.2\% for the first three components, against 47.1\% reported by
\cite{padilla2020} (Table~\ref{tab:fpca_summary}). This agreement
validates the computational pipeline and establishes that the
differences reported for the alternative bases are attributable to the
basis expansion and not to a change in the underlying data. The
exploratory analysis further confirmed the three properties that
motivate the comparison: time-varying cross-sectional volatility, with
peaks around 1982 and 2020; strong heterogeneity in country-level
volatility, with the standard deviation of Venezuela (8.76\%) more than
three times that of Guatemala (2.40\%); and crisis episodes that differ
across countries in timing, magnitude, and duration.

\subsection{Implementation of the three basis families}

B-spline, Fourier, and wavelet (Haar, DB2, Sym4) expansions were
implemented within a single Python pipeline, with the discrete wavelet
transform constructed from first principles as an explicit orthogonal
matrix. This construction guarantees exact reconstruction at zero
threshold and makes every step of the wavelet representation
transparent and reproducible. A practical consequence of this design is
that the soft-threshold parameter $\lambda$ defines a continuous family
of representations, ranging from the observed data ($\lambda = 0$) to a
smooth approximation as $\lambda$ increases, so that the degree of
smoothing becomes an explicit and interpretable control rather than an
implicit property of the basis.

\subsection{Comparison metrics at matched reconstruction quality}

Calibrating the wavelet threshold, country by country, to the
reconstruction error of the B-spline fit proved necessary for a fair
comparison: without it, the wavelet basis would trivially attain zero
error, and any difference in sparsity or FPCA structure would confound
the effect of the basis family with the effect of the smoothing level.
Under this calibration, the wavelet representations require more
nonzero coefficients than the B-spline fit (38.0 of 128 for Sym4 and
38.8 of 128 for DB2, against 19.5 knots on average). The result
therefore does not support the claim that wavelets are more economical
than B-splines in absolute coefficient count on this dataset. What the
calibration does show is that, at the same mean error, the two families
distribute that error differently. Because the B-spline error target is
proportional to each country's noise level, the coefficient count at matched
error is a relative measure and does not by itself rank countries by
regularity. The
shape of the sparsity--fidelity curves (Figure~\ref{fig:sparsity_tradeoff})
is consistent with the Besov-space characterization of spatially
inhomogeneous regularity: the less volatile economies tolerate high
compression with little increase in error, whereas the most volatile
economies do not.

\subsection{Preservation of crisis episodes and volatility structure}

At matched reconstruction error, neither the B-spline nor the wavelet
family uniformly preserves the crisis episodes. The B-spline follows large,
multi-year episodes, such as the 2015--2020 contraction of Venezuela, but
flattens isolated single-year shocks; the thresholded wavelets retain some
of those isolated shocks in attenuated form, but soft thresholding also
shrinks the large episodes, so that the Venezuelan contraction reaches about
$-19\%$ against observed values near $-30\%$
(Figure~\ref{fig:venezuela_comparison}). For Argentina, neither family
recovers the 2002 and 2020 contractions at the matched error level. The
Fourier basis is the least suited of the three to this dataset: the
power spectrum shows no well-defined business-cycle periodicity
(Figure~\ref{fig:fourier_spectrum}), and increasing the number of
retained frequencies reduces the error at the cost of spurious
oscillations over otherwise stable periods, of the type predicted by the
Gibbs phenomenon. The wavelet energy decomposition adds information that
neither global basis provides: Colombia concentrates approximately 72\%
of its energy in the coarsest approximation level, Argentina places its
largest share in the 16--32 year detail level, and Venezuela distributes it
across several levels. The time scale at which each country's variability operates
is thus an observable output of the representation and not an assumption
imposed on it.

\subsection{Effect of the basis on FPCA parsimony and interpretability}

The parsimony of the FPCA depends strongly on the basis, ranging from
three components for 80\% of the variance under B-spline smoothing to
eight under the unsmoothed data (Table~\ref{tab:fpca_summary}). This
parsimony must not be equated with quality. Under B-spline smoothing,
the first component explains 53.5\% of the variance but reduces to the
post-1980 growth slowdown, because the crisis variation has been removed
before the decomposition. Under Sym4 thresholding at $\lambda = 5$, five
components are required, and the first component (33.8\%) describes a
pattern that can be read as a regional commodity cycle: expansion in the
1970s, contraction during the 1982--1990 decade, recovery during the
2000s commodity boom, and a sharp decline in 2020. Relative to
\cite{padilla2020}, whose B-spline representation produced leading
components that each explain less than 18\% of the variance, the
thresholded wavelet representation yields fewer, more interpretable
components without discarding the crisis episodes. The threshold
$\lambda$ moves the FPCA continuously between these regimes, reaching the
B-spline parsimony at $\lambda = 12$, which confirms that the differences
observed across bases are largely differences in what each
representation removes from the data.

\subsection{General conclusion}

For volatile and multiscale economic series such as Latin American GDP
growth, a single global smoothness criterion is inadequate, and a low
reconstruction error or a high proportion of variance explained by the
first functional principal component is not, by itself, evidence of an
adequate functional representation. Among the families compared, the
wavelet basis at moderate thresholding provides the most informative
representation: it retains localized shocks, exposes country-specific
multiscale structure, and produces FPCA modes with a direct economic
reading. The B-spline basis remains adequate for the more stable
economies, whose trajectories are dominated by a smooth long-run
component. The choice of basis should therefore be justified with
respect to the regularity of the data and reported as part of the
analysis, which is precisely the step that the reviewed literature
(Section~\ref{sec:soa_gap}) and the original study of
\cite{padilla2020} leave unexamined.

%---------------------------------------------------------
\section{Limitations}
\label{sec:limitations}

The conclusions above are subject to the following limitations, which
delimit their scope.

\begin{enumerate}
    \item \textbf{Sample size.} The panel contains $N = 20$ trajectories.
    The consistency and asymptotic normality results of
    Section~\ref{sec:fpca_theory} justify FPCA as $N \to \infty$; with
    twenty countries, the estimated eigenvalues and eigenfunctions,
    particularly beyond the first two components, are subject to
    substantial sampling variability that has not been quantified in
    this thesis.

    \item \textbf{Qualitative assessment of crisis preservation.} The
    preservation or suppression of crisis episodes (specific objective 4)
    was assessed through visual comparison of reconstructions and through
    the energy decomposition, not through a dedicated quantitative metric.
    The RMSE is computed over the full domain and does not distinguish
    errors concentrated in crisis years from errors distributed over
    stable periods.

    \item \textbf{Boundary treatment.} Both the discrete Fourier
    transform and the wavelet transform assume a periodic signal. The
    wavelet transform is applied after extending each series of 65
    observations circularly to length 128, so part of the 128
    coefficients represents the periodic extension rather than observed
    data, and the discontinuity between the first and last observations
    may generate boundary artefacts in both families. The sparsity
    figures should be read with this in mind.

    \item \textbf{Smoothing parameter selection.} The B-spline smoothing
    level was set by a rule proportional to a robust estimate of each
    country's noise level, and the wavelet thresholds used for the FPCA
    ($\lambda \in \{5, 12\}$) were fixed values common to all countries.
    Neither was selected by a formal data-driven criterion such as
    generalized cross-validation for splines or the SURE and universal
    thresholds for wavelets.

    \item \textbf{Scope of the comparison.} The comparison is restricted
    to three classical, non-data-driven basis families and to a single
    variable observed at annual frequency. Adaptive-knot splines,
    data-driven orthonormal bases, and Bayesian shrinkage approaches,
    reviewed in Chapter~\ref{chap:state_of_art}, were deliberately
    excluded, and the conclusions do not extend to them.

    \item \textbf{Missing data.} Missing values for the early years,
    notably for Cuba (1960--1971) and El Salvador (1960--1965), were
    imputed by interpolation before the basis expansion. Although the
    reconstruction metrics exclude imputed points, the imputed values
    enter the FPCA and may affect the estimated components in the early
    part of the domain.
\end{enumerate}

%---------------------------------------------------------
\section{Recommendations}
\label{sec:recommendations}

\subsection{Recommendations for the application of FDA to economic series}

\begin{enumerate}
    \item The basis system should be treated as a modelling decision and
    justified with respect to the regularity of the data. Before fitting
    a smooth global basis, a preliminary multiresolution or spectral
    analysis, such as the wavelet energy decomposition and the power
    spectrum used in Chapter~\ref{chap:results}, can indicate whether the
    series are dominated by a smooth trend, by a periodic component, or by
    localized shocks.

    \item Basis families should be compared at matched reconstruction
    quality. Comparisons made at the default settings of each method
    conflate the effect of the basis family with the effect of the
    smoothing level.

    \item The proportion of variance explained by the leading functional
    principal components should be reported together with the shape of
    the corresponding eigenfunctions and with the smoothing level used,
    since a high proportion of explained variance may result from the
    removal of variation rather than from its efficient organization.

    \item For panels with marked heterogeneity in volatility across
    units, the smoothing level should be set per unit or chosen by a
    spatially adaptive method, rather than through a single global
    parameter.
\end{enumerate}

\subsection{Recommendations for future work}

\begin{enumerate}
    \item Define a quantitative crisis-preservation metric, for instance
    the reconstruction error restricted to windows around documented
    crisis years, or the detection rate of structural breaks in the
    reconstructed series, so that specific objective 4 can be evaluated
    formally.

    \item Replace the circular extension with boundary-corrected wavelets
    or symmetric extension, and evaluate the Fourier basis on detrended
    or symmetrically extended series, in order to separate the Gibbs
    oscillations due to interior structural breaks from those due to the
    boundary.

    \item Select smoothing parameters by data-driven criteria, such as
    generalized cross-validation or restricted maximum likelihood for
    penalized splines and SURE or level-dependent thresholds for
    wavelets, and include translation-invariant (stationary) wavelet
    transforms to remove the dependence of the reconstruction on the
    alignment of the dyadic grid.

    \item Strengthen the B-spline baseline with spatially adaptive
    variants, such as free-knot or locally penalized splines, to
    determine whether the advantages observed for wavelets derive from
    multiresolution itself or from local adaptivity, which splines can
    also provide.

    \item Quantify the sampling variability of the FPCA results through
    bootstrap confidence bands for the eigenvalues and eigenfunctions,
    given the small number of trajectories in the panel.

    \item Extend the analysis to quarterly data, to other macroeconomic
    variables of the region (inflation, exchange rates, terms of trade),
    and to a detailed study of the Ecuadorian economy, for which the
    literature review found no functional data applications.

    \item Use the thresholded wavelet representation as input to
    functional clustering of countries, building on the tighter grouping
    of the Central American economies observed in the score plots
    (Section~\ref{subsec:score_plots}).
\end{enumerate}
